\section{Optical transport uniform in the initial critical cutoff}
\label{sec:v81-transport}
We now use the joint volume estimate on the same survivor graph as
before. The factor $a^{-1}$ in the earlier all-incidence estimate
can be removed without changing the initial event or discarding
internal grazing scales. This improvement also applies to the
compression cost of the nongrazing boundary pullback.

The definitions of $E_h^{\varepsilon,a}$,
$E_h^{\varepsilon,\mathrm{gr},a}$, $O_h^{\varepsilon,a}$ and
$\mathcal F_h^a$ are unchanged. In particular $a$ is an initial
lower cutoff for $|g|$, and the flux keeps the first-exit selection
and each incident physical component.

\begin{theorem}[A cutoff-independent optical height bound]
\label{thm:v81-optical-height}
For each fixed $a,r>0$,
\begin{equation}\label{eq:v81-optical-height}
 \limsup_m\sup_{R,n,k}m^2
       \|\bar O_r^{\varepsilon,a}
                  +\bar E_r^{\varepsilon,\mathrm{gr},a}\|_\infty
                           \le Cr,
\end{equation}
with a constant independent of $a$ and $\varepsilon$.
For initial sources in $a\le |g|<b$, if
$b'=(b^2+2\Lambda_\circ r)^{1/2}<b_1$, the stronger estimate is
\begin{equation}\label{eq:v81-critical-annulus}
 \limsup_m\sup_{R,n,k}m^2
    \|\bar O_r^{\varepsilon,[a,b)}
          +\bar E_r^{\varepsilon,\mathrm{gr},[a,b)}\|_\infty
                      \le Cr(b')^\delta.
\end{equation}
All parameters in these statements are fixed before the collision
limit. The constants are uniform over those choices.
\end{theorem}
\begin{proof}
Repeat the negative-optical-variation argument of
Theorem~\ref{thm:v80-noncritical-optical}, retaining $|g|^{-1}$
in the target insertion rather than replacing it by $a^{-1}$.
The radial monotonicity gives $|g(\Phi_sx)|\ge a$. The common
coarea kernels of Lemma~\ref{lem:v80-common-coarea} therefore yield
\begin{equation}\label{eq:v81-drift-before-limit}
 O_h^{\varepsilon,a}(t)+E_h^{\varepsilon,\mathrm{gr},a}(t)
                            \le C\int_0^h v^a(t+s)\,ds
\end{equation}
for almost every $(t,h)$. The entire positive starting mass of a
grazing exit is paid by negative optical variation as before; a
vanishing boundary density is not used to erase that mass.
Averaging gives $2CrA_rv^a(t)$, and
\eqref{eq:v81-full-truncation} proves \eqref{eq:v81-optical-height}.

For the annulus, the radial speed additionally gives
$a\le |g(\Phi_sx)|<b'$ for $s<r$ before the first exit.
Replace $v^a$ by $v^{a,b'}$ in the comparison, and apply
\eqref{eq:v81-truncated-average}. This proves
\eqref{eq:v81-critical-annulus}. The comparison is on the full
source and its injective survivor graph, not on an estimate for
each word followed by a cardinality bound.
\end{proof}

\begin{theorem}[Nongrazing pullback with a uniform paid drift]
\label{thm:v81-boundary-height}
At almost every $(t,h)$,
\begin{equation}\label{eq:v81-boundary-before-limit}
 E_h^{\varepsilon,a}(t)+O_h^{\varepsilon,a}(t)
        \le C\mathcal F_h^a(t)+C\int_0^h v^a(t+s)\,ds.
\end{equation}
Consequently the original scalar endpoint budget improves to
\begin{align}
 \mathcal E_M\le{}&C_MB^{-1/192}
   +\mathfrak C_{\varepsilon(B),M}(a,r)+C\mathfrak F_M(a,r)
                                                       \notag\\
   &+Cr+C_\Theta(1+\varepsilon(B)^{-1})\sqrt r.
                                      \label{eq:v81-endpoint-budget}
\end{align}
The complete critical initial defect and nongrazing flux are exactly
those in \eqref{eq:v80-endpoint-budget}; only their paid drift has
changed. The leading constant no longer requires $r/a\to0$.
\end{theorem}
\begin{proof}
For a nongrazing first exit use
$\omega(0)\le\omega(\tau-)+\int_0^\tau[-X\omega]_+\,ds$,
dividing by the initial angular Jacobian. The terminal term is at
most the fixed angular-ratio constant times the assigned outgoing
flux. The integral is compared with $v^a$ using the physical
Jacobian and the same target coarea kernels. Grazing exits and
surviving optical losses are disjoint source cases and use
\eqref{eq:v81-drift-before-limit}. Thus the inverse-Jacobian
compression cost is paid, not set equal to one, and
\eqref{eq:v81-boundary-before-limit} follows.

Apply the original source-shift inequality before taking any
supremum, split its positive defect by initial $|g|<a$ or
$|g|\ge a$, and average in $h$. The second part uses the new
comparison and \eqref{eq:v81-full-truncation}; the first part is
precisely $\mathfrak C$. The already proved guard and angular
errors are unchanged. The collision limit is taken after fixing
$B,\varepsilon(B),a,r$ and every auxiliary band; the auxiliary band
is enlarged afterwards. This gives \eqref{eq:v81-endpoint-budget}.
\end{proof}

\subsection{The untruncated critical comparison}
There is also a genuine improvement without a lower critical cutoff,
but its kernel is still singular. Let $v^\eta$ denote the complete
source with insertion $I_m|g|^{-\eta}$, obtained by adding the
bounded-away-from-zero part to \eqref{eq:v81-fractional-source}.
For $0<\eta<1/2$, its fixed-width normalized height is bounded.
Since $|g(\Phi_sx)|^2\ge2\lambda_\circ s$, the same positive
variation proof gives the extended nonnegative inequality
\begin{equation}\label{eq:v81-fractional-abel}
 O_h^\varepsilon(t)+E_h^{\varepsilon,\mathrm{gr}}(t)
       \le C_\eta\int_0^h s^{-(1-\eta)/2}v^\eta(t+s)\,ds.
\end{equation}
This uses an actual untruncated mixed insertion, not a formal
integration of an uncapped spectral remainder.

Neither \eqref{eq:v81-fractional-abel} nor the uniformity in $a$
of \eqref{eq:v81-optical-height} justifies removing the initial
cutoff inside the collision-normalized limsup. At fixed count,
monotone convergence recovers the full source as $a\downarrow0$.
That fixed-count assertion is not an exchange of the two limits.
In particular the critical tail can concentrate at roofs depending
on the collision count even though all its fixed-width averages
are controlled. The outgoing nongrazing flux also remains a
surface source, not a consequence of the volume insertion theorem.

The organizing assertion remains \eqref{eq:v76-target}. To finish
it through \eqref{eq:v81-endpoint-budget}, it remains to control
the two displayed complete heights with positive outer choices
$a(B),r(B)$, while
$(1+\varepsilon(B)^{-1})\sqrt{r(B)}\to0$. No further condition
$r(B)/a(B)\to0$ is imposed by the now estimated optical drift.
Scalar closure continues to precede every path normalization and
every likelihood consequence on positive arithmetic reference classes.
